\section{Introduction}\label{sec:introduction}
Two universities with the same number of patents can produce them through
different knowledge processes: one draws on a narrow set of external sources
and refines an established core; the other connects technical elements that
rarely appear together and is later reused by organizations outside its own
field. Patent counts and citations record the output of both institutions but
not the difference between them \citep{griliches1990,hall2005,nagaoka2010}.
Understanding how institutional knowledge structures change therefore requires
reading patents as relational documents---linked to classifications, texts,
earlier references, later citing records, and dates---rather than as units to
be counted.

Information science already treats documents this way. Bibliographic coupling
and co-citation identify intellectual connections \citep{kessler1963,small1973,price1965};
knowledge maps represent domains and their change over time
\citep{borner2003,chen2006,vaneck2010}; patent overlays extend these methods to
technology \citep{rafols2010,kay2014,leydesdorff2014,yan2017}. What is missing
is a way to move from a map, which shows where an institution is located at one
moment, to a trajectory, which shows how its knowledge relations change and
whether that change exceeds sampling and measurement noise.

This paper proposes such a connection for Chinese university patents. The
central idea is simple. Patent relations of different kinds---shared
classifications, textual proximity, and citations---express different
knowledge relationships and become observable at different times. Keeping them
separate, and dating each one, makes it possible to describe an institution by
four process dimensions that are observable when an invention is published or
shortly after (acquisition of external knowledge, recombination of technical
elements, early emergence, and documented reuse), and to assess structural
transformation independently, as a sustained change in the institution's
technological composition across periods. The two research questions are:
\begin{enumerate}[label=\textbf{RQ\arabic*.}]
\item How can acquisition, recombination, emergence, and reuse be measured from
patent relations so that each retains its own meaning and its own observation
time?
\item Which university knowledge profiles and trajectories does a
university-only patent and citation corpus support, once patent volume,
technological composition, and coverage are taken into account?
\end{enumerate}

The contribution is a traceable link between knowledge organization and
longitudinal institutional analysis: every profile value and every trajectory
statement can be followed back to dated patent records and the specific
relation that supports it. The paper specifies the framework, the measures,
and the reporting design, and states explicitly which parts a university-only
corpus can and cannot support. Empirical values remain to be computed; no
finding is asserted below.

\section{Related work}\label{sec:related}
Patents are selective records of knowledge production. Not all inventions are
patented, quality varies, and legal and organizational strategies shape the
record \citep{griliches1990,nagaoka2010}. Citations require particular care:
examiner-added references weaken their reading as inventor learning
\citep{alcacer2006}, and knowledge-flow estimates depend on technological
controls and comparison groups \citep{jaffe1993,thompson2005}. Patents are
therefore treated here as documented technical objects with dated relations,
not as complete histories of invention.

Three relation types are used in patent mapping. Classification systems
organize technical categories \citep{leydesdorff2008}; text exposes
similarities that shared codes miss \citep{whalen2020,kelly2021,sarica2020};
citations record sourcing and later use. Overlay maps place organizations in a
common knowledge space \citep{rafols2010,kay2014}, and relatedness research
links positions in that space to entry and exit \citep{kogler2013,rigby2015}.
Layout movement, however, can arise from a changed sample or algorithm rather
than from changed knowledge. This study keeps visual coordinates fixed and
measures change in the underlying distributions.

The process vocabulary comes from innovation research. Absorptive capacity
concerns the recognition and use of external knowledge \citep{cohen1990};
exploration and exploitation \citep{march1991} and search depth and scope
\citep{katila2002} describe learning orientations; combinatorial and
atypical-combination studies motivate measuring unusual connections separately
from their later impact \citep{fleming2001,uzzi2013,youn2015,verhoeven2016,wang2017}.
Emergence is more than novelty: its attributes include growth, coherence,
impact, and uncertainty \citep{rotolo2015}, of which scientometric methods
operationalize a subset \citep{small2014}. First citation and later accumulation
carry different temporal information \citep{huang2019}. Multilayer network
research shows why heterogeneous relations should not be summed into one
adjacency matrix \citep{kivela2014,boccaletti2014}. The present framework draws
on these sources but makes a narrower claim: patent references index documented
sourcing, not capability; rare code pairs index unusual combination, not
usefulness; later citing records index bounded reuse, not learning.

\section{Framework: relations, time, and transformation}\label{sec:framework}
Technological knowledge evolution is defined here as change in the sources,
combinations, recognized directions, and documented reuse of an institution's
technical knowledge, together with sustained change in how that knowledge is
organized. Patents are the observable objects; universities are the contexts in
which they accumulate and connect. No patent or institution is assigned an
intrinsic developmental type.

Figure~\ref{fig:framework} shows the design. Four dimensions describe a
university-period profile,
\begin{equation}
\mathbf{x}_{ut}=(A_{ut},R_{ut},E_{ut},D_{ut})^{\mathsf T},
\label{eq:profile}
\end{equation}
where $A$ (acquisition) and $R$ (recombination) are observable when inventions
are published, $E$ (emergence) uses the current cohort against its own
historical reference, and $D$ (reuse) requires a fixed follow-up horizon.
Transformation, $T_{u,t\rightarrow t+1}$, is not a fifth input. It is a
cross-period statistic on the distribution of technical categories and pairs,
assessed separately with its own null. Keeping $T$ outside the profile prevents
a circular reading in which future structural change is used to describe a
state and the state is then said to explain that change.
\begin{figure}[tbp]
\centering
\begin{tikzpicture}[x=1cm,y=1cm]
% time axis
\draw[arr,gray] (-6.6,1.7)--(6.6,1.7) node[right,font=\scriptsize,gray]{time};
\node[note,gray] at (-4.6,2.05) {historical reference $B_t$};
\node[note,gray] at (0,2.05) {cohort $W_t$ (first publication)};
\node[note,gray] at (4.6,2.05) {follow-up $(\tau_p,\tau_p+H]$};
\draw[gray,dashed] (-2.3,1.5)--(-2.3,-4.0);
\draw[gray,dashed] (2.3,1.5)--(2.3,-4.0);
% four concurrent dimensions, placed by the time at which they become observable
\node[box,text width=3.6cm] (a) at (-4.6,0) {\textbf{Knowledge Sourcing} $A$\\External backward references: reach and diversity};
\node[box,text width=3.6cm] (r) at (0,0) {\textbf{Recombination} $R$\\Rarity of IPC pairs against $B_t$};
\node[box,text width=3.6cm] (e) at (0,-2.5) {\textbf{Technological Exploration} $E$\\Semantic novelty $\times$ field growth against $B_t$};
\node[box,text width=3.6cm] (d) at (4.6,0) {\textbf{Knowledge Diffusion} $D$\\External citing breadth and first-uptake velocity};
\node[note] at (-4.6,-1.35) {relation: citation (backward)};
\node[note] at (4.6,-1.35) {relation: citation (forward)};
\node[note] at (0,-3.75) {relation: classification and text};
% profile
\node[box,text width=12.6cm,fill=gray!8] (prof) at (0,-5.3) {\textbf{University-period profile} $\mathbf{x}_{ut}=(A,R,E,D)$\\Four concurrent dimensions, each dated by its own observation window; relation layers never summed};
\draw[arr] (a.south)--(a.south |- prof.north);
\draw[arr] (e.south)--(e.south |- prof.north);
\draw[arr] (d.south)--(d.south |- prof.north);
% dynamic transformation across periods (the fifth dimension)
\node[box,text width=12.6cm] (t) at (0,-7.7) {\textbf{Dynamic Transformation} $T_{u,t\rightarrow t+1}$ (cross-period, the fifth dimension)\\Divergence of subclass and pair distributions between $W_t$ and $W_{t+1}$ against a within-university null; sustained if $W_{t+2}$ stays closer to the new composition};
\draw[arr,dashed] (prof.south)--(t.north);
\end{tikzpicture}
\caption{Framework. Four concurrent dimensions are placed at the time they become observable and are drawn from distinct patent relations; dynamic transformation is a separate cross-period statistic---the fifth dimension---kept outside the concurrent profile to avoid circularity. No arrow denotes a causal effect or a compulsory stage.}
\label{fig:framework}
\end{figure}


Two temporal rules follow. First, each dimension is dated by the earliest
observation that determines it, and no indicator uses information published
after its own observation window. Second, a profile at one period is a state,
not a trajectory; trajectory statements require ordered profiles for the same
institution and evidence that observed change exceeds a within-university null.

Four interpretive trajectory categories are retained as \emph{hypotheses to be
checked against these observables}, not as classes to be imposed. A
\emph{maintenance} trajectory shows a stable technical core (low $T$, no
sustained coverage growth). An \emph{expansion} trajectory shows sustained
growth in rarefied subclass coverage without unusual pairing. A
\emph{recombination} trajectory shows repeatedly high $R$ supported by textual
evidence of new element connections. A \emph{transformation} trajectory
requires $T$ above its null in one transition and persistence of the new
composition in the following period. An institution may follow different
categories at different times, mixed cases remain visible, and the categories
form no performance ranking.

\section{Data and observable relations}\label{sec:data}
\subsection{Corpus and its boundaries}
The intended source is the research team's database of Chinese university
patents with their citation records. Table~\ref{tab:data} lists the coverage
items that must be audited before analysis. Its rows determine which relations
in the framework are observable and are therefore part of the design rather
than a formality.
\begin{table}[tbp]\centering\fontsize{9}{11}\selectfont
\begin{tabularx}{\textwidth}{|Y|l|Y|}
\hline
\rowcolor{black} \color{white} Coverage item & \color{white} Value & \color{white} Required definition or audit\\\hline
Database/provider; extraction date & \TBD & Frozen source and query manifest\\\hline
Retrieval dates; analytic dates & \TBD & First-publication cohorts; mature end date\\\hline
Roster institutions; linked institutions & \TBD & Dated roster and name crosswalk\\\hline
Raw publication records & \TBD & Before publication/application deduplication\\\hline
Eligible invention units & \TBD & Families, or applications if families unavailable\\\hline
Eligible university-periods & \TBD & Three-year non-overlapping cohorts\\\hline
IPC subclasses; main groups & \TBD & Version-harmonized vocabulary sizes\\\hline
Usable title/abstract records & \TBD & Count and coverage percentage\\\hline
Backward-citation links & \TBD & Resolved earlier public invention units\\\hline
Forward-citation links within $H$ & \TBD & Unique family pairs, scope declared\\\hline
External citing organizations & \TBD & University-only or broader coverage\\\hline
Excluded units by reason & \TBD & Mutually exclusive first-failure counts\\\hline
\end{tabularx}
\caption{Dataset scale and coverage, to be computed from the frozen source database. No value or date range in this table is an empirical estimate.}
\label{tab:data}
\end{table}


The primary cohort comprises published Chinese invention applications with at
least one university applicant at first publication, regardless of later grant
status; utility models and designs are excluded. Institutional coverage is
defined by a dated university roster, with historical names, mergers, and
affiliated hospitals resolved through time-specific linkage; ambiguous matches
remain auditable rather than silently assigned. Publication and grant versions
of an application are collapsed. If family identifiers are available, the
invention unit is a simple family represented by its earliest eligible
document; otherwise it is the deduplicated application, and the unit is
declared. A patent with $m_p$ university applicants contributes
$\omega_{up}=1/m_p$ to each; full counts are a sensitivity.

Three properties of a university-only corpus constrain the design and are
stated here rather than discovered later. (i)~Backward references from
university patents are observable, but the cited documents' own attributes
are available only when the cited record is itself in the corpus or its
classification is carried in the reference field. (ii)~Forward citations to
university patents are observable; the citing documents' applicants are
observable only where the citation record carries applicant fields, and the
citing documents' \emph{own} forward links are generally not available. Reuse
depth beyond one citation step is therefore excluded from the primary measure.
(iii)~Chinese invention patents receive few forward citations in their first
year, so any indicator conditioned on early recognition will be zero for most
records. Emergence is consequently measured in an early-information form that
does not depend on recognition, with a recognized variant as a secondary
measure.

\subsection{Observation time}
First publication is the main clock. $W_t=[b_t,b_{t+1})$ is a non-overlapping
three-year cohort and $B_t=[b_t-3,b_t)$ its historical reference; reuse uses a
horizon of $H=5$ years. Mature profiles require $b_{t+1}+H\leq d_{\mathrm{extract}}$,
and cohorts without a full historical reference are excluded. The number of
mature cohorts this leaves is a property of the database, to be reported in
Table~\ref{tab:data}; if fewer than three consecutive mature cohorts exist,
trajectory analysis is restricted to description and the transformation test
is reported as infeasible rather than approximated.

A university-period enters profile analysis with at least 20 attributed
inventions and at least 80\% usable attributed mass for the fields each
dimension requires. Confirmed absence of citations is zero; an unavailable
field is missing and never enters a denominator. For dimension $j$ the eligible
subset is $P^j_{ut}$ with mass $M^j_{ut}=\sum_{p\in P^j_{ut}}\omega_{up}$,
reported as a share of total mass.

\subsection{Three relation layers}
Figure~\ref{fig:pipeline} shows how the corpus is decomposed. The technology
layer has IPC main-group nodes linked by fractional co-occurrence,
\begin{equation}
w_{ij}^{(t)}=\sum_{p\in W_t:\,m_p^C\geq2}
\frac{\ind\{i,j\in\mathcal{C}_p\}}{\binom{m_p^C}{2}},\quad i<j,
\label{eq:cooc}
\end{equation}
where $\mathcal{C}_p$ is the set of distinct main groups on patent $p$; each
multi-coded invention contributes one unit of pair mass. The semantic layer
has patent nodes linked by cosine similarity of character 2--4-gram TF--IDF
vectors of normalized titles and abstracts, with vocabulary and IDF fitted on
$B_t$ only (features in at least five and at most 80\% of reference documents,
capped at 100{,}000; log TF, smoothed IDF, $L_2$ normalization). The citation
layer has patent nodes with directed edges from cited to citing document,
duplicate family pairs collapsed, within-family links removed, and citing
publication required to be strictly later. Applicant-overlap self-citations are
excluded from the primary reuse graph. A patent--IPC incidence relation
connects the layers; their adjacency matrices are never summed. Display graphs
(normalized edges, top-10 neighbor union; mutual 15-nearest-neighbor semantic
graph; Leiden communities, \citealp{traag2019}) use fixed coordinates across
periods; all indicators use unpruned weights.
\begin{figure}[tbp]\centering
\begin{tikzpicture}[x=1cm,y=1cm]
\node[box,text width=12.3cm] (source) at (0,0) {\textbf{University-only patent and citation corpus}\\Dated university roster; deduplicated invention units; coverage audit (Table~\ref{tab:data})};
\node[box,text width=3.25cm,minimum height=2.3cm] (ipc) at (-4.5,-3.0) {\textbf{Technology layer}\\IPC main-group pairs\\Fixed layout; $B_t$ reference};
\node[box,text width=3.25cm,minimum height=2.3cm] (sem) at (0,-3.0) {\textbf{Semantic layer}\\Title--abstract TF--IDF\\Novelty against $B_t$};
\node[box,text width=3.25cm,minimum height=2.3cm] (cit) at (4.5,-3.0) {\textbf{Citation layer}\\Backward: sourcing\\Forward: one-step reuse};
\foreach \n in {ipc,sem,cit}{\draw[arr] (source.south)--++(0,-.35)-|(\n.north);}
\node[box,text width=12.3cm] (prof) at (0,-5.9) {\textbf{University-period profiles} $A$, $R$, $E$, $D$ \quad and \quad \textbf{cross-period} $T$\\Denominators, coverage shares, and observation horizons retained};
\foreach \n in {ipc,sem,cit}{\draw[arr] (\n.south)--++(0,-.35)-|(prof.north);}
\node[box,text width=3.8cm,minimum height=2.2cm] (land) at (-4.3,-8.9) {\textbf{Landscape}\\Fixed-layout map;\\period overlays};
\node[box,text width=3.8cm,minimum height=2.2cm] (dist) at (0,-8.9) {\textbf{Profiles}\\Distributions; size and field adjustment};
\node[box,text width=3.8cm,minimum height=2.2cm] (traj) at (4.3,-8.9) {\textbf{Trajectories}\\Ordered profiles; $T$ events vs.\ null; sequence tests};
\foreach \n in {land,dist,traj}{\draw[arr] (prof.south)--++(0,-.35)-|(\n.north);}
\node[box,text width=12.3cm] (valid) at (0,-11.6) {\textbf{Cases and sensitivity}\\Blinded case reading; representation, unit, window, and threshold sensitivity};
\foreach \n in {land,dist,traj}{\draw[arr] (\n.south)--++(0,-.3)-|(valid.north);}
\end{tikzpicture}
\caption{Data-to-inference workflow. Three relation layers are kept separate up to the profile; transformation is computed across periods from the technology layer alone.}
\label{fig:pipeline}
\end{figure}


\section{Measures}\label{sec:measures}
\subsection{Acquisition}
Let $q_{ut}$ be the attributed share of valid inventions with at least one
resolved external backward reference, and let $\pi_{ut,k}$ be the fractional
distribution of those references over a fixed, version-harmonized IPC subclass
vocabulary of size $K_t$. Then
\begin{equation}
A_{ut}=q_{ut}\left[\tfrac12+\tfrac12\,
\frac{-\sum_k\pi_{ut,k}\log\pi_{ut,k}}{\log K_t}\right],
\label{eq:acquisition}
\end{equation}
so that confirmed absence of external sourcing gives $A=0$ and the positive
baseline distinguishes narrow sourcing from none. Both components are
reported. $A$ measures documented source reach and diversity
\citep{shannon1948,stirling2007}, not absorptive capacity \citep{cohen1990}.

\subsection{Recombination}
Let $c^-_{ij,t}$ be the historical fractional pair mass of main groups
$(i,j)$ in $B_t$, $Q_t$ the number of possible pairs in the pre-established
vocabulary, and $C_t=\sum_{i<j}c^-_{ij,t}>0$. With $\alpha=1/2$,
\begin{align}
p^-_{ij,t}&=\frac{c^-_{ij,t}+\alpha}{C_t+\alpha Q_t},&
r_{ij,t}&=\frac{-\log p^-_{ij,t}}{-\log\{\alpha/(C_t+\alpha Q_t)\}},\\
R_{ut}&=\frac{1}{M^R_{ut}}\sum_{p\in P^R_{ut}}\omega_{up}\,
\frac{\sum_{(i,j)\in\mathcal{L}_p}r_{ij,t}}{|\mathcal{L}_p|},
\label{eq:recombination}
\end{align}
where $\mathcal{L}_p$ is the set of valid mapped pairs on patent $p$.
Unresolved codes map to one recorded out-of-vocabulary category; a patent with
no valid pair is missing, not zero. Rarity indexes unusual combination, not
novelty or usefulness \citep{fleming2001,youn2015,verhoeven2016}. First-seen
pair rates and degree-preserving permutations of the patent--IPC incidence
provide complementary checks.

\subsection{Emergence}
Semantic novelty is $n_{p,t}=1-\max_{q\in B_t}s_{pq}$ over the full historical
index, excluding same-family comparators. For subclass $k$, with fractional
shares $\rho_{k,t}$ in the current cohort and $\rho^-_{k,t}$ in the reference,
\begin{equation}
g_{k,t}=\max\!\left\{0,\frac{\rho_{k,t}-\rho^-_{k,t}}
{\rho_{k,t}+\rho^-_{k,t}+10^{-12}}\right\},\qquad
E_{ut}=\frac{1}{M^E_{ut}}\sum_{p\in P^E_{ut}}\omega_{up}\,n_{p,t}\,g_{p,t},
\label{eq:emergence}
\end{equation}
where $g_{p,t}$ averages over the patent's subclasses. $E$ is thus a mean of
patent-level novelty--growth products: novelty and field growth must coincide
in the same invention. It is an early-information measure available at
publication. A recognized variant $E^{\mathrm{rec}}$ multiplies each product by
an indicator of external citation within one year; given constraint~(iii) in
Section~\ref{sec:data}, it is reported with its nonzero share as a secondary
measure and is not the primary emergence indicator. Neither variant covers the
full emergence construct \citep{rotolo2015,small2014}; technical coherence is
assessed through cases.

\subsection{Reuse}
Within $(\tau_p,\tau_p+H]$, let $b_p$ count distinct external citing
organizations where applicant fields are available and distinct external
citing families otherwise (the level is recorded), and let
$v_p=1-\Delta_p/H$ when the first external citation arrives after $\Delta_p$
years, with $v_p=0$ for uncited patents. Each component is transformed to its
midrank among positive values in the same primary subclass and cohort
($\phi(0)=0$; back-off to IPC section and then cohort below 30 positives), and
\begin{equation}
D_{ut}=\frac{1}{2M^D_{ut}}\sum_{p\in P^D_{ut}}\omega_{up}
\left[\phi^b(b_p)+\phi^v(v_p)\right].
\label{eq:diffusion}
\end{equation}
$D$ measures \emph{documented reuse among observed citers}. Because the
citing side of a university-only corpus is incomplete, $D$ does not measure
economy-wide diffusion, and a two-step depth component is not part of the
primary measure; it is computed as a sensitivity only where descendants'
forward links exist. First uptake and accumulated breadth carry different
information \citep{huang2019} and are also reported separately.

\subsection{Transformation}
Let $\mathbf{a}_{ut}$ and $\mathbf{r}_{ut}$ be the university's normalized
subclass and pair distributions on shared vocabularies. Then
\begin{equation}
T_{u,t\rightarrow t+1}=\frac{
\JSD(\mathbf{a}_{ut},\mathbf{a}_{u,t+1})+
\JSD(\mathbf{r}_{ut},\mathbf{r}_{u,t+1})}{2\log2}\in[0,1],
\label{eq:transformation}
\end{equation}
with $\JSD$ the Jensen--Shannon divergence \citep{lin1991}. A transformation
event requires $T$ above the 95th percentile of a within-university null that
pools the two adjacent cohorts and permutes their period labels while
preserving cohort sizes, confirmed under size-matched subsampling. A
\emph{sustained} event additionally requires the third period to be closer to
the new composition than to the old and still distinguishable from the old.
Mergers and IPC revision events are audited before any event is retained.
$T$ measures institutional reorganization of a portfolio, not a technological
paradigm shift \citep{dosi1982,funk2017}; diversity is a distinct property
\citep{stirling2007,porter2009}.

\section{Analysis and reporting design}\label{sec:analysis}
The analysis proceeds in three steps, each answerable from the corpus described
above.

\emph{Landscape.} The technology layer is laid out once on the pooled corpus
(Figure~\ref{fig:map}A); period overlays reuse the same coordinates
(Figure~\ref{fig:map}B) so that entry, persistence, and exit of technical
groups reflect the data rather than the layout. Coverage growth is evaluated
against rarefied subclass coverage and portfolio size, since a larger sample
alone makes more fields visible. Disagreement between the semantic and
classification views marks cases for inspection.

\emph{Profiles.} Distributions, zero shares, and denominators of $A,R,E,D$ are
reported for eligible university-periods (Table~\ref{tab:descriptive}), together
with pairwise correlations and their association with log patent volume, field
composition, and period. Residual variation after these adjustments is the
evidence that relational profiles add information beyond size and field.
Profiles are then displayed on their original scales with whole-university
bootstrap intervals \citep{efron1979}, and universities are placed in profile
space (Figure~\ref{fig:states}A). Model-based grouping of profiles (e.g., a
Gaussian mixture with stability checks, \citealp{fraley2002,hennig2007}) is a
possible extension once the number of mature university-periods is known; it
is not required for the questions posed here and is not asserted in advance.

\emph{Trajectories.} For universities with consecutive eligible cohorts,
ordered profiles are displayed alongside $T$ and its null
(Figure~\ref{fig:states}B). The four trajectory categories of
Section~\ref{sec:framework} are checked against their stated observable
criteria; the counts of supported, mixed, and unsupported cases are reported.
The proposition that sustained recombination precedes transformation is tested
by comparing the observed frequency of ``high-$R$ followed by sustained $T$
event'' with its frequency under within-university order permutations that
preserve each university's profile values. A nonpositive or unstable contrast
leaves the proposition unsupported.

\emph{Cases and sensitivity.} Twelve trajectories sampled across profile
regions and size strata are assessed by two technically qualified raters from
dated patent evidence without model values; agreement, disagreements, and
adjudication are reported \citep{cohen1960}. Sensitivity covers invention
unit, attribution, text representation (TF--IDF versus a frozen
Chinese-capable encoder, \citealp{reimers2019}), window length, citation
horizon and self-citation scope, eligibility thresholds, and the $E$ and $D$
variants (Table~\ref{tab:robustness}). Common-sample agreement is reported
separately from sample retention. All secondary associations are exploratory.

\section{Results}\label{sec:results}
\emph{Evidence status.} Values are not yet computed; the paragraphs below fix
what will be reported and contain no observed finding.

\subsection{Coverage and landscape}
The analytic population comprises \pending{universities}, \pending{unique
inventions}, and \pending{university-periods} in \pending{mature cohorts}.
Table~\ref{tab:data} separates raw retrieval, deduplication, and eligibility,
with the leading exclusion reasons \pending{reasons and counts}. Of the citation
layer, \pending{share} of forward links carry citing applicant fields, which
fixes the level at which $b_p$ is computed. Figure~\ref{fig:map} reports
\pending{node, edge, and community counts} and the retained mass; entry,
persistence, and exit of technical groups are \pending{observed patterns}.
\begin{figure}[tbp]\centering
\figpanel{0.48\textwidth}{A. Common technological space}{
Data required: IPC nodes and weighted pairs.\\[.6em]
Position: one fixed reference layout.\\
Node area: fractional invention mass.\\
Edges: normalized co-occurrence.\\[.6em]
\textbf{Empirical map not yet computed.}}
\hfill
\figpanel{0.48\textwidth}{B. Period overlays}{
Data required: period-specific IPC mass.\\[.6em]
Reuse panel A coordinates.\\
Show entry, persistence, and exit.\\
Report retained mass and omitted nodes.\\[.6em]
\textbf{Empirical overlays not yet computed.}}
\caption{Technological knowledge space: analysis specification awaiting source data. The panels contain no synthetic nodes, invented clusters, or empirical coordinates.}
\label{fig:map}
\end{figure}


\subsection{Profiles}
Table~\ref{tab:descriptive} reports each dimension's distribution and zero
share. Correlations among $A,R,E,D$ are \pending{coefficients and intervals},
and the share of variance remaining after adjustment for volume, field, and
period is \pending{estimate}. Figure~\ref{fig:states}A shows where universities
fall in profile space and which regions are populated; $E^{\mathrm{rec}}$ is
nonzero for \pending{share} of university-periods.
\begin{table}[tbp]\centering\fontsize{9}{11}\selectfont
\begin{tabularx}{\textwidth}{|Y|c|c|c|c|c|c|}
\hline
\rowcolor{black} \color{white} Measure & \color{white} Valid $n$ & \color{white} Mean & \color{white} SD & \color{white} Median & \color{white} IQR & \color{white} Zero \%\\\hline
Knowledge sourcing ($A$) & \TBD & \TBD & \TBD & \TBD & \TBD & \TBD\\\hline
Recombination ($R$) & \TBD & \TBD & \TBD & \TBD & \TBD & \TBD\\\hline
Technological exploration ($E$) & \TBD & \TBD & \TBD & \TBD & \TBD & \TBD\\\hline
Recognized exploration ($E^{\mathrm{rec}}$, secondary) & \TBD & \TBD & \TBD & \TBD & \TBD & \TBD\\\hline
Knowledge diffusion ($D$) & \TBD & \TBD & \TBD & \TBD & \TBD & \TBD\\\hline
Dynamic transformation ($T$, per transition) & \TBD & \TBD & \TBD & \TBD & \TBD & ---\\\hline
\end{tabularx}
\caption{Distribution of the five dimensions across eligible university-periods and of $T$ across consecutive eligible transitions. Component summaries, missing counts, and coverage shares accompany each row in the output file.}
\label{tab:descriptive}
\end{table}


\subsection{Trajectories and transformation}
Among \pending{count} universities with consecutive eligible cohorts,
\pending{count} show at least one $T$ event above its null and \pending{count}
show a sustained event; \pending{count} events are removed by the merger and
IPC-revision audit. Trajectory categories supported by their observable
criteria number \pending{counts by category, mixed, unsupported}. The contrast
between observed and permuted ``recombination then transformation'' sequences
is \pending{effect and interval}. Figure~\ref{fig:states}B displays sampled
ordered profiles with their $T$ values.
\begin{figure}[tbp]\centering
\figpanel{0.48\textwidth}{A. Profile space}{
Data required: $A,R,E,D$ per eligible university-period.\\[.6em]
Original scales; pairwise panels or a\\
secondary PCA view for orientation.\\
Point size: attributed invention mass.\\
Whole-university bootstrap intervals.\\[.6em]
\textbf{No profile values computed.}}
\hfill
\figpanel{0.48\textwidth}{B. Ordered trajectories}{
Data required: consecutive cohorts and $T$ with null.\\[.6em]
Horizontal axis: publication cohort.\\
Rows: prespecified sampled universities.\\
Bars: $A,R,E,D$; markers: $T$ above null,\\
sustained events, audited mergers.\\[.6em]
\textbf{No trajectories assigned.}}
\caption{Profile and trajectory displays: analysis specification. Trajectory categories are read from the stated observable criteria and the independent $T$ test, not from a fitted clustering.}
\label{fig:states}
\end{figure}


\subsection{Cases and sensitivity}
Rater agreement on \pending{count} cases is \pending{agreement and kappa};
principal disagreements concern \pending{themes}. Table~\ref{tab:robustness}
reports the settings that change profile values or trajectory counts most;
dependence on a single representation, on sparse citations, or on one
eligibility threshold bounds the interpretation.
\begin{table}[tbp]\centering\fontsize{9}{11}\selectfont
\begin{tabularx}{\textwidth}{|P{3.1cm}|Y|P{2.4cm}|}
\hline
\rowcolor{black} \color{white} Test family & \color{white} Alternative and comparison & \color{white} Observed result\\\hline
Invention unit & Application versus family; common-sample profile change & \TBD\\\hline
Attribution & University-fractional versus all-applicant and full counts & \TBD\\\hline
Text representation & Character TF--IDF versus frozen Chinese-capable encoder; novelty rank agreement & \TBD\\\hline
Window length & Non-overlapping 2-, 3-, and 4-year cohorts; pattern-category agreement & \TBD\\\hline
Diffusion horizon and scope & $H=3$ versus 5; self-citation exclusion; organization- versus family-level $b_p$ & \TBD\\\hline
Eligibility & Minimum 10/20/30 inventions; completeness 70/80/90\% & \TBD\\\hline
Exploration variant & $E$ versus $E^{\mathrm{rec}}$; nonzero shares & \TBD\\\hline
Transformation null & Label permutation versus size-matched subsampling; event counts & \TBD\\\hline
Cases & Blinded rater agreement; cases rejecting the category reading & \TBD\\\hline
\end{tabularx}
\caption{Sensitivity and case results, pending computation. Each row reports the common-sample size, the change in profile values or pattern counts, an interval where applicable, and the change in coverage.}
\label{tab:robustness}
\end{table}

\FloatBarrier

\section{Discussion}\label{sec:discussion}
The framework's contribution to information science is representational.
Each inference is tied to the relation that supports it: an unusual IPC pair
supports a combination claim; semantic distance from a dated reference corpus
supports a relative novelty claim; a later external citing record supports a
bounded reuse claim; divergence between successive category distributions,
against a within-university null, supports a reorganization claim. Retaining
these origins makes agreement and disagreement among representations visible
and connects patent mapping \citep{kay2014,leydesdorff2014,yan2017} to
longitudinal institutional analysis without collapsing the relations into one
score.

The design also states what a university-only corpus can support. It can
describe sourcing, combination, early emergence, and reuse among observed
citers, and it can test whether a portfolio's composition changed more than
its own year-to-year variation. It cannot establish economy-wide diffusion,
reuse chains beyond one step, or the causes of change. If recurring profiles or
ordered sequences do not exceed their nulls, the framework remains a
descriptive organization of evidence rather than an account of developmental
mechanisms, and that outcome is reported as such.

Other limits persist. Patents omit scientific, tacit, and unpatented knowledge.
Citations do not observe learning or contractual transfer \citep{alcacer2006}.
Novelty is relative to a corpus and a representation. Institutional events
such as mergers, recruitment, and policy affect both profiles and later change.
Profiles should therefore be read with their coverage and uncertainty,
consistent with responsible use of research metrics \citep{hicks2015}, and not
as a ranking.

\section{Conclusion}\label{sec:conclusion}
This paper connects the contents of patent records, their dated relations, and
the time at which each relation becomes observable into a framework for
describing how university knowledge structures change. Four process dimensions
are measured from separate relation layers; transformation is assessed
independently as sustained compositional change. The framework is specified
against the actual boundaries of a university-only patent and citation corpus,
so that its empirical claims---when computed---are limited to what the data
observe. Extensions with enterprise citing records would allow reuse to be
traced beyond the academic system, and publication--patent links would separate
scientific from technological sourcing.

\section*{Data and code availability}
The source database and computed outputs were not available for this revision;
database permissions and any deposit must be verified by the authors. The
accompanying LaTeX package contains the manuscript and display sources only.
Provenance, output fields, and audit requirements follow traceable data
stewardship principles \citep{wilkinson2016}.
